\subsection{正根}

设 C 是 R 的一个室，且 \(\mathrm{B}(\mathrm{C})=\{\alpha_{1},\ldots,\alpha_{l}\}\) 是相应的 R 的基。由 C 在 V（相应地，在 \(V^{*}\)）上定义的序关系，是与 V（相应地，与 \(V^{*}\)）的向量空间结构相容的序关系，并且其中满足 \(\geqslant0\) 的元素是 \(\alpha_{i}\)（相应地，\(\alpha_{i}^{-}\)）的线性组合，且系数为 \(\geqslant0\)。在这些序关系之一中为正的元素称为关于 C 正，或关于基 \(\mathrm{B}(\mathrm{C})\) 正。这些序关系也由 \(C^{-}\) 定义，正如将 V 与 \(V^{*}\) 等同并使用在 \(\mathrm{W}(\mathrm{R})\) 下不变的内积所看到的。根据 Th. 2, no. 5，\(V^{*}\) 中的元素是 \(\geqslant0\) 当且仅当其在 C 上的取值为 \(\geqslant0\)。V 中的元素 x 是 \(\geqslant0\) 当且仅当它在 \(C^{-}\) 上的取值为 \(\geqslant0\)，或者等价地，当且仅当 \((x|y)\geqslant0\) 对所有 \(y\in C\) 成立。

\(\overline{C}\) 的元素是 \(\geqslant 0\)（相对于 \(C\)），由 Chap. V, § 3, no. 5 的 Lemma 6 可知。但是，\(\geqslant 0\) 的元素集合（相对于 \(C\)）一般不同于 \(\overline{C}\)（cf.~Plate X, Systems \(A_2\) 、\(B_2\) 、\(G_2\)）。

定理 3. 每个根都是 B(C) 中元素的同号整数系数组合。特别地，每个根关于 C 不是正的就是负的。

若 \(\alpha \in \mathbb{R}\)，则 \(L_{\alpha}\) 是 \(\alpha\) 的核，不与 \(C^{\cdot}\) 相交，因此 \(\alpha\) 要么为 \(>0\)（在整个 \(C^{\cdot}\) 上），要么为 \(< 0\)（在整个 \(C^{\cdot}\) 上），从而得到第二个断言。还需证明 \(\alpha\) 属于子群 \(P\)，它是 \(V\) 中由 \(B(C)\) 生成的子群；可以假设 \(\alpha\) 不可分。现在群 P 显然在 \(s_{\gamma}\) 下保持不变，其中 \(\gamma \in \mathrm{B}(\mathrm{C})\)，因而根据 Th. 2 也在 \(\mathrm{W}(\mathrm{R})\) 下保持不变。由于 \(\alpha\) 具有 \(w(\beta)\) 的形式，其中 \(w \in \mathrm{W}(\mathrm{R})\) 且 \(\beta \in \mathrm{B}(\mathrm{C})\)（cf.~命题 15），所以 \(\alpha \in P\)。证毕。

记 \(R_{+}(C)\) 为关于 C 为正的根的集合。于是，

\[
R = R _ {+} (C) \cup (- R _ {+} (C))
\]

是 \(R\) 的一个划分。

推论。设 \(\gamma\) 是根的整数系数组合，\(\alpha\) 是不可分根。若 \(\gamma\) 与 \(\alpha\) 成比例，则 \(\gamma \in Z\alpha\)。

由第 5 号的命题 15，可以选择 C 使得 \(\alpha \in B(C)\)。由定理 3，

\[
\gamma = \sum_ {\beta \in B (C)} n _ {\beta} \beta \text {with} n _ {\beta} \in \mathbf {Z}.
\]

因此，若 \(\gamma\) 与 \(\alpha\) 成比例，则 \(\gamma = n_{\alpha}\alpha\)，这就证明了该推论。

现在令 S 为反射 \(s_{\alpha}\) 的集合，其中 \(\alpha \in \mathrm{B}(\mathrm{C})\)，并令 T 为 S 在 W 下的共轭元素的并集。对于 \(\alpha \in \mathrm{B}(\mathrm{C})\) 和 \(w \in W\)，T 中的元素 \(t = ws_{\alpha}w^{-1}\) 是正交反射 \(s_{\beta}\)，其中根 \(\beta = w(\alpha)\)；反过来，对于任意不可分根 \(\beta\)，存在 \(w \in W\)，使得 \(\alpha = w^{-1}(\beta) \in \mathrm{B}(\mathrm{C})\)（命题 15），并且 \(s_{\beta} = ws_{\alpha}w^{-1} \in T\)。因此，得到一个双射 \(\psi\)，从不可分根集合到 \(\{\pm1\} \times T\)，它把不可分根 \(\beta\) 对应到二元组 \((\varepsilon, s_{\beta})\)，其中 \(\varepsilon = +1\) 当 \(\beta\) 为正，且 \(\varepsilon = -1\) 当 \(\beta\) 为负。

另一方面，(W, S) 是一个 Coxeter 系统（定理 2），可以应用 Chap. IV, § 1, no. 4 的结果。我们已经看到，若 w 是 W 中长度（相对于 S）为 q 的元素，则存在一个含有 q 个元素的 T 的子集 \(T_{w}\)，使得当 \(w = s_{1} \ldots s_{q}\)，其中 \(s_{i} \in S\)，并且

\[
t _ {i} = s _ {1} \dots s _ {i - 1} s _ {i} s _ {i - 1} \dots s _ {1}
\]

（对于 \(1 \leqslant i \leqslant q\)），则 \(\mathrm{T}_w = \{t_1, \ldots, t_q\}\)。回顾一下，我们还在第 1 节第 4 号定义了一个数 \(\eta(w, t)\)，当 \(w \in \mathrm{W}\) 且 \(t \in \mathrm{T}\) 时，其值为 \(+1\)（当 \(t \notin \mathrm{T}_w\) 时），为 \(-1\)（当 \(t \in \mathrm{T}_w\) 时）。最后，回顾一下，如果我们用下式定义映射 \(\mathrm{U}_w\)，从集合 \(\{\pm 1\} \times \mathrm{T}\) 到自身，

\[
\mathbf {U} _ {w} (\varepsilon , t) = (\varepsilon \eta (w ^ {- 1}, t), w t w ^ {- 1}),
\]

映射 \(w \mapsto \mathrm{U}_w\) 是从 \(W\) 到集合 \(\{\pm 1\} \times T\) 的置换群的同态（Chap. IV, §1, no. 4, Lemma 1）。

命题 17. 假设 R 是约化的，并令 \(w \in W\)、\(\alpha \in R\)。

\begin{enumerate}
\def\labelenumi{(\roman{enumi})}
\item
有 \(\psi(w(\alpha)) = U_w(\psi(\alpha)).\)
\item
  假设 \(\alpha\) 为正。根 \(w(\alpha)\) 为负，当且仅当
\end{enumerate}

\[
\eta (w ^ {- 1}, s _ {\alpha}) = - 1,
\]

换言之，当 \(s_{\alpha} \in \mathrm{T}_{w^{-1}}\) 时。

\begin{enumerate}
\def\labelenumi{(\roman{enumi})}
\setcounter{enumi}{2}
\tightlist
\item
当且仅当 \(\eta(w,s_{\alpha})=-1\)，室 C 和 \(w(C)\) 位于超平面 \(L_{\alpha}\) 的两侧。换言之，集合 \(T_{w}\) 由分隔 C 和 \(w(C)\) 的各壁对应的反射组成。
\end{enumerate}

设 \(\beta \in \mathrm{B}(\mathrm{C})\)，并置 \(s = s_{\beta}\)。显然 \(\mathrm{T}_s = \{s\}\)，从而

\[
\mathrm{U} _ {s} (\varepsilon , t) = \left\{ \begin{array}{l l} (\varepsilon , s t s ^ {- 1}) & \text { if } t \neq s \\ (- \varepsilon , s) & \text { if } t = s. \end{array} \right.\tag{12}
\]

另一方面，设 \(\rho = \sum_{\gamma \in B(C)} n_{\gamma}(\rho) \gamma\) 是一个正根。置

\[
s (\rho) = \sum_ {\gamma \in B (C)} n _ {\gamma} (s (\rho)) \gamma .
\]

若 \(\rho \neq \beta\)，则存在 \(\gamma \in \mathrm{B}(\mathbf{C})\)，满足 \(\gamma \neq \beta\)，使得 \(n_{\gamma}(\rho) > 0\)，并且有 \(n_{\gamma}(s(\rho)) = n_{\gamma}(\rho) > 0\)（第 1 号的公式 (5)）。因此 \(s(\rho)\) 是正的。立即得到

\[
\psi (s (\varepsilon . \rho)) = \left\{ \begin{array}{l l} (\varepsilon , s s _ {\rho} s ^ {- 1}) & \text {if} \rho \neq \beta \\ (- \varepsilon , s) & \text {if} \rho = \beta . \end{array} \right.\tag{13}
\]

比较 (12) 和 (13) 现在表明 \(\mathrm{U}_{s}(\psi(\gamma)) = \psi(s(\gamma))\) 对所有根 \(\gamma\) 及所有 \(s \in S\) 成立。由于 S 生成 W，(i) 得证。

另一方面，称 \(w(\alpha)\) 为负，等价于

\[
\psi (w (\alpha)) = (- 1, w s _ {\alpha} w ^ {- 1}),
\]

或者根据 (i)，等价于 \(\mathrm{U}_w(\psi (\alpha)) = (-1, ws_{\alpha}w^{-1})\)。如果此外 \(\alpha\) 为正，则 \(\psi (\alpha) = (+1,s_{\alpha})\)，且 \(\mathrm{U}_w(\psi (\alpha)) = (\eta (w^{-1},s_{\alpha}),ws_{\alpha}w^{-1})\)，故得到 (ii)。

最后，根据 (ii)，\(\eta(w,s_{\alpha})=-1\) 当且仅当根 \(\alpha\) 和 \(w^{-1}(\alpha)\) 中一个为正而另一个为负。这等价于 \((\alpha|x).(w^{-1}(\alpha)|x)=(\alpha|x).( \alpha|w(x))<0\) 对所有 \(x\in C\) 成立，从而得到 (iii) 的第一个断言。(iii) 的第二个断言立即可得。

推论 1. 设 \(\beta \in \mathrm{B}(\mathbb{C})\)。反射 \(s_{\beta}\) 将与 \(\beta\) 不成比例的正根进行置换。

我们立即归结到 R 约化的情形。在这种情况下，我们的断言由 (ii) 以及 \(T_{s_{\beta}} = s_{\beta}\) 这一事实得到。

推论 2. 假设 R 是约化的。设 \(w \in W\)，设 \(q\) 是 \(w\) 关于 S 的长度（Chap. IV, § 1, no. 1），并设 \(w = s_1 \ldots s_q\) 是 \(w\) 的一个约化分解。设 \(\alpha_1, \ldots, \alpha_q\) 是 B(C) 中与 \(s_1, \ldots, s_q\) 对应的元素。置

\[
\theta_ {i} = s _ {q} s _ {q - 1} \dots s _ {i + 1} (\alpha_ {i}), \quad i = 1, \dots , q.
\]

根 \(\theta_{i}\) 都 \(>0\)、彼此不同，\(w(\theta_{i}) < 0\)，并且每个满足 \(\alpha > 0\) 且 \(w(\alpha) < 0\) 的根都等于某个 \(\theta_{i}\)。

令 \(X\) 为满足 \(\alpha > 0\) 且 \(w(\alpha) < 0\) 的根的集合。由 (ii)，

\[
\operatorname{Card} (X) = \operatorname{Card} \left(T _ {w ^ {- 1}}\right) = l \left(w ^ {- 1}\right) = l (w) = q.
\]

另一方面，若 \(\alpha \in X\)，显然存在 \(i \in [1, q]\)，使得

\[
s _ {i + 1} \ldots s _ {q} (\alpha) > 0 \quad \text {and} \quad s _ {i} s _ {i + 1} \ldots s _ {q} (\alpha) <   0.
\]

由 Cor. 1，这意味着 \(s_i s_{i+1} \ldots s_q(\alpha) = \alpha_i\)，因而 \(\alpha = \theta_i\)。于是集合 X 包含于 \(\theta_i\) 的集合中。由于 \(\text{Card}(X) = q\)，这只有在 X 等于 \(\theta_i\) 的集合且这些元素彼此不同时才可能。因此推论得证。

推论 3. 假设 R 是约化的。W 中存在唯一的最长元素 \(w_0\)。它的长度等于正根的数目，并且 \(w_0\) 将室 C 变为 -C。我们有 \(w_0^2 = 1\)，且 \(l(ww_0) = l(w_0) - l(w)\) 对所有 \(w \in W\) 成立。

显然 -C 是一个室。因此存在 W 中的元素 \(w_{0}\) 将 C 变为 -C。于是 \(w_{0}(\alpha) < 0\) 对所有正根 \(\alpha\) 成立，推论 3 的前两个断言立即由推论 2 得到。我们有 \(w_{0}^{2}(C) = C\)，所以 \(w_{0}^{2} = 1\)。最后，若 \(w \in W\)，则由命题 17 (iii)，长度 \(l(w)\)（相应地，\(l(ww_{0})\)）等于分隔 C 与 \(w(C)\)（相应地，\(ww_{0}(C) = -w(C)\)）的壁数。由于 \(w(C)\) 和 \(-w(C)\) 位于每一面墙的两侧，\(l(w) + l(w_{0})\) 等于墙的总数，也就是 \(l(w_{0})\)。

命题 18. 设 \(x \in V\)。以下三个性质等价：

\begin{enumerate}
\def\labelenumi{(\roman{enumi})}
\item
  \(x \in \overline{C};\)
\item
  \(x \geqslant s_{\alpha}(x)\) 对所有 \(\alpha \in \mathrm{B}(\mathrm{C})\) 成立（相对于 \(\mathrm{C}\) 定义的序关系）；
\item
   \(x \geqslant w(x)\) 对所有 \(w \in W\) 成立。
\end{enumerate}

由于 \(s_{\alpha}(x) = x - \langle x, \alpha' \rangle \alpha\)，且 \(\overline{\mathbf{C}}\) 是由 \(x \in V\) 构成的集合，满足 \(\langle x, \alpha' \rangle \geqslant 0\) 对所有 \(\alpha \in \mathrm{B}(\mathbf{C})\) 成立，(i) 与 (ii) 的等价性显然。另一方面，显然 (iii) \(\Longrightarrow\) (ii)。我们证明 (i) \(\Longrightarrow\) (iii)。设 \(x \in \overline{\mathbf{C}}\)，且 \(w \in W\)。我们对长度 \(l(w)\) 的 \(w\) 归纳。\(l(w) = 0\) 的情形是平凡的。若 \(l(w) \geqslant 1\)，则 \(w\) 可写成 \(w = w's_{\alpha}\) 的形式，其中 \(\alpha \in \mathrm{B}(\mathbf{C})\) 且 \(l(w') = l(w) - 1\)。于是

\[
x - w (x) = x - w ^ {\prime} (x) + w ^ {\prime} \left(x - s _ {\alpha} (x)\right).
\]

归纳假设表明 \(x - w'(x)\) 是正的。另一方面，

\[
w ^ {\prime} (x - s _ {\alpha} (x)) = w \left(s _ {\alpha} (x) - x\right) = - \langle x, \alpha^ {*} \rangle w (\alpha).
\]

现在 \(s_{\alpha} \in T_{w^{-1}}\)，命题 17 (ii) 表明 \(w(\alpha) < 0\)。故结论成立。

推论。元素 \(x \in \mathbf{C}\) 当且仅当有 \(x > w(x)\) 对所有 \(w \in \mathbf{W}\) 且 \(w \neq 1\) 成立。

命题 19. 设 \((\beta_{i})_{1\leqslant i\leqslant n}\) 是关于室 C 的正根序列，且 \(\beta_{1}+\beta_{2}+\cdots+\beta_{n}\) 是一个根。则存在置换 \(\pi\in S_{n}\)，使得对所有 \(i\in\{1,2,\ldots,n\}\)，\(\beta_{\pi(1)}+\beta_{\pi(2)}+\cdots+\beta_{\pi(i)}\) 都是根。

我们对 n 归纳，\(n \leqslant 2\) 时命题显然。置 \(\beta = \beta_{1} + \cdots + \beta_{n}\)。于是 \(\sum_{i=1}^{n} (\beta | \beta_{i}) = (\beta | \beta) > 0\)，所以存在一个指标 k，使得 \((\beta | \beta_{k}) > 0\)。若 \(\beta = \beta_{k}\)，则 n = 1，因为对所有 i 都有 \(\beta_{i} > 0\)。否则 \(\beta - \beta_{k}\) 是一个根（第 3 号的 Th. 1 的 Cor.）；于是只需将归纳假设应用于 \(\beta - \beta_{k} = \sum_{i \neq k} \beta_{i}\)。

推论 1. 设 \(\alpha \in \mathrm{R}_{+}(\mathrm{C})\)。则 \(\alpha \in \mathrm{B}(\mathrm{C})\)，当且仅当 \(\alpha\) 是两个正根之和。

若 \(\alpha\) 是两个正根之和，则定理 3 表明 \(\alpha \in \mathrm{B}(\mathrm{C})\)。若 \(\alpha \notin \mathrm{B}(\mathrm{C})\)，则定理 3 表明 \(\alpha = \sum_{k=1}^{n} \beta_k\)，其中 \(\beta_k \in \mathrm{B}(\mathrm{C})\) 对所有 \(k\) 成立，且 \(n \geqslant 2\)。必要时置换 \(\beta_k\)，可假设 \(\sum_{k=1}^{n-1}\) 是一个根（命题 19），因而 \(\alpha\) 是正根 \(\sum_{k=1}^{n-1} \beta_k\) 与 \(\beta_n\) 之和。

推论 2. 设 \(\varphi\) 是从 \(R\) 到交换群 \(\Gamma\) 的映射，具有以下性质：

\begin{enumerate}
\def\labelenumi{\arabic{enumi})}
\item
  \(\varphi(-\alpha) = -\varphi(\alpha)\) 对 \(\alpha \in \mathbf{R}\) 成立；
\item
  若 \(\alpha \in \mathrm{R}\)、\(\beta \in \mathrm{R}\) 且 \(\alpha + \beta \in \mathrm{R}\)，则 \(\varphi(\alpha + \beta) = \varphi(\alpha) + \varphi(\beta)\)。
\end{enumerate}

令 \(\mathbf{Q}\) 为 \(\mathbf{V}\) 中由 \(\mathbf{R}\) 生成的子群。于是 \(\varphi\) 延拓为从 \(\mathbf{Q}\) 到 \(\Gamma\) 的同态。

设 B 是 R 的一个基。设 \(\psi\) 是从 Q 到 \(\Gamma\) 的唯一同态，并且在 B 上与 \(\varphi\) 一致。只需证明 \(\psi(\alpha) = \varphi(\alpha)\)，其中 \(\alpha\) 是相对于 B 的正根。我们有 \(\alpha = \beta_{1} + \cdots + \beta_{m}\)，其中对所有 i 都有 \(\beta_{i} \in B\)，且对所有 h 都有 \(\beta_{1} + \cdots + \beta_{h} \in R\)（命题 19）。我们对 m 归纳证明 \(\psi(\alpha) = \varphi(\alpha)\)。当 m = 1 时显然成立。归纳假设给出

\[
\psi \left(\beta_ {1} + \dots + \beta_ {m - 1}\right) = \varphi \left(\beta_ {1} + \dots + \beta_ {m - 1}\right),
\]

且我们有 \(\psi (\beta_m) = \varphi (\beta_m)\)，所以 \(\psi (\alpha) = \varphi (\alpha)\)，这就证明了该推论。

 对于满足 \(\alpha = \sum_{\beta \in \mathrm{B}(\mathbb{C})} n_{\beta} \beta\) 且属于 \(\mathbb{R}\) 的任意根，记 \(\Upsilon(\alpha)\) 为满足 \(\beta \in \mathrm{B}(\mathbb{C})\) 且 \(n_{\beta} \neq 0\) 的元素的集合。此外，注意到 \(\mathrm{B}(\mathbb{C})\) 可以与由 \(\mathrm{W}(\mathbb{R})\) 和 \(s_{\alpha_i}\) 构成的 Coxeter 系统的图的顶点集等同（cf.~Chap. IV, § 1, no. 9 and Chap. V, § 3, no. 2）。

推论 3. a) 设 \(\alpha \in \mathbb{R}\)。则 \(\mathrm{Y}(\alpha)\) 是 \(\mathrm{B}(\mathrm{C})\) 的连通子集（Chap. IV, Appendix）。

\begin{enumerate}
\def\labelenumi{\alph{enumi})}
\setcounter{enumi}{1}
\tightlist
\item
  设 Y 是 B(C) 的非空连通子集。则 \(\sum_{\beta\in Y}\beta\) 属于 R。
\end{enumerate}

为了证明 \(a\) )，可假设 \(\alpha\) 为正。我们对 \(\operatorname{Card}(\mathbf{Y}(\alpha))\) 归纳；当 \(\operatorname{Card}(\mathbf{Y}(\alpha)) = 1\) 时断言显然。由命题 19，存在 \(\beta \in \mathrm{B}(\mathrm{C})\)，使得 \(\alpha - \beta \in \mathrm{R}\)。令 \(p\) 为最大的满足 \(\geqslant 0\) 且 \(\gamma = \alpha - p\beta \in \mathrm{R}\) 的整数。由于 \(\gamma - \beta \notin \mathrm{R}\) 且 \(\gamma + p\beta \in \mathrm{R}\)，有 \((\gamma|\beta) \neq 0\)（命题 9）；因此 \(\beta\) 至少与 \(\mathbf{Y}(\gamma)\) 中的一个元素相连。但 \(\mathbf{Y}(\alpha) = \mathbf{Y}(\gamma) \cup \{\beta\}\)，且 \(\mathbf{Y}(\gamma)\) 根据归纳假设是连通的。因此 \(\mathbf{Y}(\alpha)\) 是连通的，这就证明了 \(a\) )。

现在设 Y 是 B(C) 的非空连通子集；我们对 \(\operatorname{Card}(Y)\) 归纳证明 \(\sum_{\beta\in Y}\beta\) 是一个根。\(\operatorname{Card}(Y)\leqslant1\) 的情形是平凡的。假设 \(\operatorname{Card}(Y)\geqslant2\)。由于 X 是一片森林（Chap. V, § 4, no. 8, Prop. 8），Y 是一棵树，并且有一个终端顶点 \(\beta\)（Chap. IV, Appendix）。集合 \(Y-\{\beta\}\) 是连通的，且其中一个元素与 \(\beta\) 相连。由归纳假设，\(\alpha=\sum_{\gamma\in Y-\{\beta\}}\gamma\in R\)，又由于 \((\alpha|\beta)<0\)，可知 \(\alpha+\beta\in R\)（Th. 1）。证毕。

